% Paper-ready replacement subsection. Round-12 development endpoints only.
% Sources, all matched comparisons, and scope are documented in
% docs/INTERVENTIONAL_DISCRIMINATION_EVIDENCE.md.
\subsection{Interventional discrimination}
\label{sec:experiment_identifiability}

Task specifications identify the inputs that may be manipulated and the
mechanisms a model must represent. The resulting dynamical model supports an
input intervention by replacing the supplied input schedule while holding its
equations and fitted parameters fixed, then simulating the response from the
available initial observations. Future target measurements are never supplied
during this rollout; predictions remain conditional on any auxiliaries permitted
by the task. We therefore evaluate free-rollout NMSE on validation conditions
absent from training, including new input amplitudes, timings, and frequencies.
This tests response prediction under changed excitation, beyond fitting the
recorded training trajectories.

The existing benchmark evidence illustrates why training error alone is
insufficient, although it does not establish a uniform benefit from task
specification. Across six cells from three system families, the prediction-only
ablation currently covers two cells and two seeds each. On the alien-device
cell at seed~1, prediction-only achieves lower training NMSE than the full method
($0.472$ versus $0.773$), but higher validation NMSE ($0.655$ versus $0.627$).
Thus a $39\%$ training-error reduction does not translate into a better response
under the held-out validation schedules. The reversal is modest: the full method
has lower validation error in only two of the four matched runs. This ablation
removes scientific requirements from both proposal generation and eligibility
checks, measuring their joint contribution.

A complementary example concerns model memory. On CSTR at seed~0, removing
persistent latent states improves training NMSE from $0.132$ to $0.123$, but
increases validation NMSE from $0.0140$ to $0.0320$ ($2.28\times$). This comparison
changes the model class rather than the specification, and the other CSTR seed
does not show the reversal. These examples motivate assessing response behavior
separately from training fit; they do not demonstrate unique causal recovery or
a large, systematic specification advantage. Since validation also selects
models, these are development comparisons, not independent test estimates.
